\documentclass[lualatex,paper=a4paper]{jlreq}
\usepackage{amsmath}
\usepackage{amssymb}
\usepackage{amsthm}
\usepackage{amsfonts}  %以上必須パッケージ
\usepackage{mathtools} %参照していない数式に番号をつけない
\usepackage{graphicx}  %図を入れる
\usepackage{multirow}  %表の列結合
\usepackage{hyperref}  %文章内にリンクを貼る
\usepackage{diffcoeff} %diff、diffpで微分記号を簡潔に書く
\usepackage{comment}   %comment環境
\usepackage{pdfpages}  %PDFページを直接挿入
\usepackage{float}     %[H]指定子でその位置に強制的に図表を配置する
\usepackage[version=4]{mhchem} %\ce{$n$H3O+}のようにして簡潔に化学式を書く
\usepackage[separate-uncertainty]{siunitx}  %SI単位を\unitと\qtyを使ってより簡便に書く
\usepackage[math-style=ISO,warnings-off={mathtools-colon,mathtools-overbracket}]{unicode-math} %Unicode文字を数式に書けるようにする
\usepackage{newunicodechar} %Unicode文字を命令化できるようにする
\usepackage{zref-clever}    %\crefで参照の種類に応じて「式」「図」「表」などと自動的に出力してくれる
\usepackage{listings}       %ソースコードを簡単に書けるようにする
%%%%%%%%%%%%%%%%%%%%%%%%%%%%%%%%%%%%%%%%%%%%%%%%%
\mathtoolsset{showonlyrefs,showmanualtags} %参照されていない数式に番号をつけない
\newunicodechar{、}{，}  %、を，に自動置換
\newunicodechar{。}{．}  %。を．に自動置換
\NewDocumentCommand\、{}{\string、} %\、で読点をそのまま出力
\NewDocumentCommand\。{}{\string。} %\。で句点をそのまま出力
\NewDocumentCommand\degC{}{\ensuremath{^\circ\symup{C}}} %摂氏温度記号
\NewDocumentCommand\abs{m}{\left|#1\right|} %絶対値の記号
\renewcommand{\lstlistingname}{コード} %listingsの環境名を「コード」に変更
%%%%%%%%%%%%%%%%%%%%%%%%%%%%%%%%%%%%%%%%%%%%%%%%%
\setmathfont{NewComputerModernMath}[
    CharacterVariant = {1} %∅(空集合)の記号を一般的なものに変更
]
\setmathfont{NewComputerModernMath}[
    range = {scr, bfscr},  %花文字のフォントを変更する、これがないと普通の筆記体として表示される
    StylisticSet = 1
]
% zref-cleverの日本語設定
\zcDeclareLanguage[allcaps]{japanese}
\zcLanguageSetup{japanese}{
  namesep   = {\nobreak} ,
  pairsep   = {と} ,        % 複数個の参照をつなぐときの接続詞
  listsep   = {、} ,        % 複数個の参照をつなぐときの接続詞（3個以上）
  lastsep   = {、および} ,  % 複数個の参照をつなぐときの最後の接続詞（3個以上）
  tpairsep  = {および} ,    % 複数個の異なる種類の参照をつなぐときの接続詞（2個）
  tlistsep  = {、} ,        % 複数個の異なる種類の参照をつなぐときの接続詞（3個以上）
  tlastsep  = {、および} ,  % 複数個の異なる種類の参照をつなぐときの最後の接続詞（3個以上）
  notesep   = {：} ,        % 参照の後に注釈をつけるときの区切り
  rangesep  = {から} ,      % 範囲指定のときの区切り
  type = section ,          % 節の引用に関する設定
    Name-sg = 第 ,
    Name-pl = 第 ,
    refbounds = {,,節,} ,
  type = subsection ,       % 小節の引用に関する設定
    Name-sg = 第 ,
    Name-pl = 第 ,
    refbounds = {,,小節,} ,
  type = subsubsection ,    % 小小節の引用に関する設定
    Name-sg = 第 ,
    Name-pl = 第 ,
    refbounds = {,,小小節,} ,
  type = paragraph ,
    Name-sg = 第 ,
    Name-pl = 第 ,
    refbounds = {,,段落,} ,
  type = appendix ,
    Name-sg = 付録 ,
    Name-pl = 付録 ,
  type = page ,             % ページの場合
    Name-sg = p. ,
    Name-pl = pp. ,
    namesep = 〜 ,
    rangesep = \textendash ,
    rangetopair = false ,
  type = figure ,
    Name-sg = 図 ,
    Name-pl = 図 ,
  type = table ,
    Name-sg = 表 ,
    Name-pl = 表 ,
  type = item ,
    Name-sg = 項目 ,
    Name-pl = 項目 ,
  type = footnote ,
    Name-sg = 脚注 ,
    Name-pl = 脚注 ,
  type = endnote ,
    Name-sg = 後注 ,
    Name-pl = 後注 ,
  type = note ,
    Name-sg = 注記 ,
    Name-pl = 注記 ,
  type = equation ,
    Name-sg = 式 ,
    Name-pl = 式 ,
    refbounds = {,(,),} ,
  type = listing ,
    Name-sg = コード ,
    Name-pl = コード ,
}
\zcsetup{ lang = japanese, nameinlink = true } % 日本語設定を有効化
%%%%%%%%%%%%%%%%%%%%%%%%%%%%%%%%%%%%%%%%%%%%%%%%%
\lstset{%
  basicstyle={\ttfamily}, % コード全体の書体
  identifierstyle={\normalsize}, % キーワード以外のコードの書体
  keywordstyle={\small\bfseries\color[rgb]{0.000, 0.000, 1.000}}, % キーワード(ifなど)の書体
  backgroundcolor={\color[gray]{0.98}}, % 背景色
  stringstyle={\small\ttfamily\color[rgb]{0.639, 0.082, 0.082}}, %文字列の書体
  commentstyle={\small\ttfamily\color[rgb]{0.000, 0.502, 0.000}}, %コメントの書体
  frame={tb}, % コードの上下に線を引く。lrを足すと左右にも線が引かれる
  breaklines=true, % 行が長くなったときに自動で改行する
  showstringspaces=false, % 文字列中のスペースを可視化(␣で表記)しない
  columns=fixed, % 等角で文字を配置する
  basewidth=0.5em, % 等幅フォントの文字幅を0.5emに設定
  numbers=left, % 行番号を左に表示
  xrightmargin=0\zw, %右側の余白を全角幅0に設定
  xleftmargin=3\zw, %左側の余白(行番号表示部分)を全角幅3に設定
  numberstyle={\scriptsize}, % 行番号の書体
  stepnumber=1, % 行番号を1行ごとに表示
  numbersep=1\zw, % 行番号とコードの間のスペースを全角幅1に設定
  lineskip=-0.5ex % 行間を詰める
}
%%%%%%%%%%%%%%%%%%%%%%%%%%%%%%%%%%%%%%%%%%%%%%%%%
% jlreqの設定を変更するためのコマンド
% 詳細: https://tug.org/texlive//Contents/live/texmf-dist/doc/latex/jlreq/jlreq-ja.html
\jlreqsetup{
    appendix_counter={ %付録に入ったときの処理
        section={
            value=0,
            the={\Alph{section}} %付番を戻してアルファベット(A, B, C...)に変更
        },
        table={
            value=0,
            the={\Alph{section}\arabic{table}} %表番号を付録のものに変更
        },
        figure={
            value=0,
            the={\Alph{section}\arabic{figure}} %図番号を付録のものに変更
        }
    },
    appendix_heading={
        section={
            font=\normalsize\bfseries,
            label_format={付録\thesection:} %付録の表記を変更
        },
        subsection={
            font=\normalsize\bfseries
        },
        subsubsection={
            font=\normalsize\bfseries
        }
    },
    caption_font=\mcfamily, %図表の題を明朝体で出力する
    caption_label_font=\mcfamily, %図表番号を明朝体にする
    caption_label_format={#1:}, %「図1:」のようにする
    caption_after_label_space=0.5\zw %:の後に全角幅半分のスペースを空ける
}
\ModifyHeading{section}{ %見出しの表示を変更
    font=\normalsize\bfseries, %見出しのサイズを変更
    lines=2, %見出しの行取りを変更
    after_lines=0
}
\ModifyHeading{subsection}{ %小見出しの表示を変更
    font=\normalsize\bfseries,
    lines=2,
    after_lines=0
}
\ModifyHeading{subsubsection}{ %小小見出しの表示を変更
    font=\normalsize\bfseries,
    lines=2,
    after_lines=0
}

\title{\LaTeX テンプレート} % タイトルを入れる
\author{電気通信大学 IX類 \\ %プログラムはもしあれば\\の直前に書く
9999999 無敵太郎} % 学籍番号と名前を入れる
\date{2112年9月3日作成 %\\ \today 更新% 再提出の際には\\の直前の%を取り除く
}


\begin{comment}
このテンプレートはhydrogen(https://hydrogen1.dev)によって作成されました。改変や再頒布、どのように扱っても問題ありません。
元々のテンプレートはこちら(https://gitlab.mma.club.uec.ac.jp/hydrogen/latex-template)にあります。
邪魔であれば、このコメントの部分(\begin{comment}からend{comment}まで)を消してください。

また、このテンプレートはLuaLaTeXを対象に作られています。
改変していないテンプレートから正常にPDFを生成できない場合、以下のURLからTeXLiveをインストールしてください。
https://mirror.ctan.org/systems/texlive/tlnet/install-tl-windows.exe
\end{comment}

\begin{document}
%%% --- ここから書き始める

%表紙が不要な場合は下の行頭の%を取り除く
%\begin{comment}
%\begin{titlepage}
%    \includepdf{cover_sample.pdf} %表紙のPDFファイルを挿入する。
%\end{titlepage}
%\end{comment}
%表紙が不要な場合は上の行頭の%を取り除く

\maketitle %タイトルの出力


\subsection{図の挿入}

図は\verb|\includegraphics|を用いると挿入できる。大きすぎるなら縮尺を変更して収めることも可能だ。\zcref{fig-sample}に見本を示しておこう。
%%
%%\begin{figure}[ht]
%%    \centering
%%    \includegraphics[scale=0.9]{sample.png}
%%    \caption{図のサンプル}
%%    \label{fig-sample}
%%\end{figure}
%%

\subsection{プログラムのソースコード}

\begin{lstlisting}[language=Python,caption={サンプルプログラム},label=code-sample]
import re
from decimal import Decimal
from fractions import Fraction

class Input_matrix:
    def __init__(self):
        self.A = [[0]]
        self.B = [[0]]
        self.C = [[0]]
        self.n = 0

    def tableinput(self,a):
        name = a
        row = 0
        row_count =0
        i = 0
        j = 0
        pre = []
        print(f"行列 {name} を入力\n入力例:1 1.2 2/3 351\n endと入力して行入力終了")

        while row != "end":
            row = input(str(i+1)+ "行目:")
            if row == "end":
                break
            pre.insert(i, row)
            row_count = len(pre[0].split())

            if re.search(r"[^\d./ -]", row) and row != "end":
                print("数字と-./以外入れるなボケ")
                del pre[i]
                continue

            elif row_count != len(pre[int(i)].split()):
                print("列数がちげーよボケ")
                del pre[i]
                continue

            print(*pre, sep="\n")
            i = i + 1
            
        a = [[Fraction(x) for x in s.split()] for s in pre]
        print(f"\n行列{name}の入力終了")
        print(*a, sep="\n")
        return a

    def numberinput(self,n):
        while True:
            n = input("倍数n:")
            if re.search(r"[^\d./ -]", n):
                print("数字と./-以外入れるなボケ")
                continue
            break
        return n


class Tablecalculate:
    def __init__(self):
        self.A = [[0]]
        self.B = [[0]]
        self.C = [[0]]
        self.n = 0
        
    def simplify(self,a):
        i = 0
        j = 0
        k = 0
        process_row = 0
        stop = False

        gyoretsushiki = False
        m = 1

        while process_row < len(a):
            i = process_row
            j= 0
        #Step1
            while a[i][j] == 0:
                if i == len(a)-1 and a[i][j] == 0:
                    i = process_row
                    j = j + 1
                    if j > len(a[0]) - 1:
                        stop = True
                        break
                    continue
                if stop == True:
                    break
                i = i + 1
            if j > len(a[0])-1:
                break
            print(f"a{i+1,j+1}に注目")
            tmp_i = i
            tmp_j = j
        #Step2
            k = process_row
            while a[k][tmp_j] == 0:
                k = k + 1
            if process_row+1 != tmp_i+1:
                for j in range(len(a[0])):
                    tmp = a[process_row][j]
                    a[process_row][j] = a[tmp_i][j]
                    a[tmp_i][j] = tmp
                print(f"{process_row+1}行目 ⇔ {tmp_i+1}行目")
                if gyoretsushiki == True:
                    m = -m
                    print(f"倍数m:{m}")
                tmp_i = process_row
                self.output(a)
        #Step3
            if a[tmp_i][tmp_j] != 1:
                k = a[tmp_i][tmp_j]
                print(f"{tmp_i + 1}行目/ {k}")
                if gyoretsushiki == True:
                    m = m * k
                    print(f"倍数m:{m}")
                for j in range(len(a[0])):
                    a[tmp_i][j] = Fraction(a[tmp_i][j] ,k)

            self.output(a)
        #Step4
            for i in range(process_row if gyoretsushiki == True else 0, len(a)):
                if a[i][tmp_j] != 0 and i != tmp_i:
                    k = a[i][tmp_j]
                    for j in range(len(a[0])):
                        a[i][j] = a[i][j] - a[tmp_i][j] * k
                    print(f"{i + 1}行目 - {tmp_i + 1}行目 × {k}")

            self.output(a)
            process_row = process_row + 1
            print("row=",process_row)
            if gyoretsushiki == True:
                print(f"倍数m:{m}")
        return a

    def output(self,a):
        b = [[0 for j in range(len(a[0]))] for i in range(len(a))]
        for j in range(len(a[0])):
            for i in range(len(a)):
                b[i][j] = str(a[i][j])
        print(" ", *b, sep="\n")
        

InputMatrix = Input_matrix()

A = InputMatrix.tableinput("A")
TableCalculate = Tablecalculate()
result = TableCalculate.simplify(A)
print("計算結果:")
print(*result, sep="\n")
\end{lstlisting}



%%% --- ここまで
\end{document}